\subsubsection*{Reading the law of cosines from a triangle}

The symbols in the law of cosines refer to the following Euclidean triangle.
The angle $\alpha$ is at the vertex $O$, the side opposite this angle is $AB$,
and the two sides adjacent to the angle are $OA$ and $OB$.

\begin{center}
\begin{tikzpicture}[scale=1.05]
    \coordinate (O) at (0,0);
    \coordinate (A) at (4.4,0);
    \coordinate (B) at (1.55,3.0);

    \fill[blue!6] (O) -- (1.0,0)
        arc[start angle=0,end angle=62.7,radius=1.0] -- cycle;

    \draw[subset] (O) -- (A)
        node[midway,below=5pt] {$d_E(O,A)$};
    \draw[subset] (O) -- (B)
        node[midway,left=7pt] {$d_E(O,B)$};
    \draw[subset] (A) -- (B)
        node[midway,above right=3pt] {$d_E(A,B)$};

    \draw[subset] (0.78,0)
        arc[start angle=0,end angle=62.7,radius=0.78];
    \node[subset label] at (0.88,0.42) {$\alpha$};

    \fill (O) circle (2.1pt) node[below left] {$O$};
    \fill (A) circle (2.1pt) node[below right] {$A$};
    \fill (B) circle (2.1pt) node[above] {$B$};

    \node[align=left,anchor=west] at (5.2,2.35) {
        adjacent to $\alpha$:\\
        $d_E(O,A)$ and $d_E(O,B)$
    };
    \node[align=left,anchor=west] at (5.2,0.65) {
        opposite $\alpha$:\\
        $d_E(A,B)$
    };
\end{tikzpicture}
\end{center}

For this labelled triangle, the classical Euclidean law of cosines is
\[
\boxed{
 d_E(A,B)^2
 =
 d_E(O,A)^2+d_E(O,B)^2
 -2d_E(O,A)d_E(O,B)
 \cos\bigl(m_{\mathrm{rad}}(\alpha)\bigr)
}.
\]

The diagram shows why the term containing the cosine uses the two sides
adjacent to $\alpha$, while the squared length on the left is the side opposite
$\alpha$. Cartesian coordinates are used only to calculate these three side
lengths.